\chapter{Runtime and Backends}

The runtime layer is intentionally thin. The compiler owns parsing, resolution,
analysis, lowering, and validation. The runtime selects a backend, maps user
inputs to DefIds, delegates execution, and maps DefId outputs back to user names.
This division keeps language semantics in compiler passes and keeps backends
focused on execution.

\section{CLI and Python API}

\pathcode{src/einlang/__main__.py} implements the command-line entry point. It
can run a file, inline source through \texttt{-c}, or standard input through
\texttt{-}. It supports backend selection, project root and stdlib overrides,
IR dumping, cProfile output, vectorization debugging, recurrence-block
debugging, reduction profiling, function and block profiling, and autodiff
debug logs.

\pathcode{src/einlang/run.py} provides a single-call Python API:
compile a source string or file, then execute it with optional named inputs. It
is the small embedding surface used by users who want to stay in Python while
trying \einlang{} programs.

\section{Execution Environment}

\pathcode{src/einlang/runtime/environment.py} implements a DefId-keyed scope
stack. Function calls and blocks push scopes; exiting a scope clears its values
so large arrays can be reclaimed. Lookup walks from the innermost scope outward.
There is no name-based runtime lookup after name resolution. This invariant
matches the compiler identity model and prevents shadowing bugs in execution.

\section{Runtime Delegation}

\pathcode{src/einlang/runtime/runtime.py} defines \texttt{EinlangRuntime}. It
constructs either a NumPy backend or IREE backend, creates a fresh backend per
execution to avoid state leaks, resolves named inputs to DefIds, resolves
\texttt{main} when present, delegates execution, and returns an
\texttt{ExecutionResult}. It also records vectorization counts from the backend
for tests and profiling.

Inputs are name-keyed at the public API boundary but DefId-keyed inside the
runtime. Before execution, the runtime resolves each input name against
top-level bindings and functions in the compiled program. Outputs travel the
opposite direction: the backend records values by DefId, and the runtime maps
those DefIds back to user-facing names in the final result.

\subsection{Runtime Delegation Algorithm}

\begin{lstlisting}[style=pythoncode,caption={Public execution boundary.}]
EinlangRuntime.execute(compilation_result, inputs):
    program = compilation_result.ir
    tcx = compilation_result.tcx
    backend = fresh_backend(selected_backend)

    input_by_defid = {}
    for name, value in user_inputs.items():
        did = resolve_top_level_input_name(name, program, tcx.resolver)
        input_by_defid[did] = value

    main_defid = resolve_main_if_present(program, tcx.resolver)
    raw_result = backend.execute(program,
                                 input_by_defid=input_by_defid,
                                 main_defid=main_defid,
                                 resolver=tcx.resolver,
                                 tcx=tcx)

    if raw_result has DefId-keyed outputs:
        return map output DefIds back to source names
    return raw_result
\end{lstlisting}

The fresh-backend step avoids state leakage between executions. The backend is
allowed to keep per-run caches such as vectorization counters and IREE modules,
but public execution starts from a clean environment.

The algorithm also marks the public/private boundary of identity. Users pass
inputs by source names because that is the ergonomic API, but the runtime
immediately resolves those names to DefIds using the compiler's resolver. From
that point onward, execution is DefId-keyed. The final mapping back to source
names is presentation only. This prevents local shadowing, imported aliases, or
generated bindings from turning into runtime name collisions.

Creating a fresh backend per run is a simple but important isolation rule.
Backends are allowed to accumulate execution counters, circular-buffer state,
compiled IREE modules, and profiling data while a program runs. Reusing that
state across calls would make tests order-dependent and could keep large arrays
alive longer than intended. The runtime therefore treats backend construction
as part of the execution transaction.

The \texttt{main} branch also keeps application-style and notebook-style uses
separate without changing compilation. When \texttt{main} exists, execution is
the value of that entry function. When it does not, top-level bindings become
observable results. Both modes still run through the same DefId environment and
backend visitor, so examples can remain script-like while larger programs can
provide an explicit entry point. The runtime boundary decides only how to start
and package execution, not how the language is analyzed.

Resolving inputs before the backend starts is important for error reporting.
An unknown public input name is a runtime API error at the boundary, whereas a
missing DefId inside backend execution is an implementation error. The
algorithm makes that distinction before any tensor allocation or compiled
backend setup occurs. This keeps user mistakes close to the public interface and
keeps backend code free to assume that input bindings already match compiler
identity.

\textit{Concrete example.} A Python caller may run a compiled program with
\texttt{\{"A": A, "B": B\}} as inputs. The runtime resolves the public names
\texttt{A} and \texttt{B} to the DefIds assigned during compilation, then passes
\texttt{\{did(A): A, did(B): B\}} to the backend. If the program produces
\texttt{C}, the backend stores the value by \texttt{did(C)}, and the runtime
maps that DefId back to the user-facing key \texttt{"C"} in the result. A
local variable also named \texttt{A} inside a helper function never enters this
public input map.

\section{Backend Interface}

\pathcode{src/einlang/backends/base.py} defines the backend contract. A backend
must execute a program, execute an individual expression, and optionally produce
target code. Backends are expected to trust compiler IR and use attached
metadata rather than re-running analysis. This is why passes such as
\texttt{LoweredExecutionFacts} matter: they put backend-facing decisions on the
IR before execution.

\section{NumPy Backend}

The NumPy backend is the main executable backend. It is split across mixins:
\pathcode{numpy_core.py} handles program execution, environment setup, function
calls, profiling, output storage, builtins, and inputs;
\pathcode{numpy_expressions.py} and related support files evaluate ordinary
expressions; \pathcode{numpy_einstein.py} and its mixins execute lowered
Einstein, reduction, comprehension, recurrence, and vectorization paths.

The backend registers fixed builtins in the environment, loads functions and
modules, evaluates statements in order, stores outputs by DefId, materializes
circular recurrence buffers only when necessary, and returns an execution
result. It also tracks vectorized, scalar, hybrid, and call-scalar counts for
Einstein clauses.

\subsection{NumPy Program Execution Algorithm}

\begin{lstlisting}[style=pythoncode,caption={NumPy backend top-level execution.}]
numpy_execute(program, input_by_defid, main_defid):
    set entry-file process variables for relative resource loading
    env = ExecutionEnvironment()
    register fixed builtins in env
    register program functions and module functions in env
    register builtin DefIds from resolver
    for did, value in input_by_defid:
        env.set_value(did, value)

    if main_defid exists:
        return call_function(env.get_value(main_defid), [])

    outputs = {}
    with env.scope():
        for statement in program.statements:
            value = statement.accept(backend_visitor)
            did, binding = output_defid_for(statement)
            if did is not None:
                store_output_value(outputs, did, value, binding.name)

        for did, value in env.current_scope_items():
            if did not in outputs:
                if value is circular recurrence buffer and must be observable:
                    value = value.materialize()
                outputs[did] = value
    return ExecutionResult(outputs=outputs)
\end{lstlisting}

Top-level execution has two modes. If a \texttt{main} function is present, the
program behaves like an entry-point application and the backend returns that
call result. Otherwise, the backend executes top-level statements in order and
collects observable bindings. Both modes share the same environment and DefId
lookup rules, so examples, tests, and embedded calls exercise the same runtime
machinery.

The output collection logic is more subtle than a dictionary copy. A top-level
statement may produce a value directly, bind a function, or create a circular
recurrence buffer whose observable value depends on later materialization
rules. The backend records outputs by DefId, checks whether ring buffers must
be materialized, and returns an \texttt{ExecutionResult} that the public runtime
can translate back to names. This preserves the compiler's identity model all
the way to the result boundary.

Registering functions and builtins before evaluating statements is the runtime
equivalent of frontend hoisting. A top-level statement may call a function that
appears later in source order, and builtins such as \texttt{shape} or
\texttt{\_\_basis\_tensor} must be available to generated derivative IR. The
backend therefore prepares the environment in layers: fixed builtins, program
functions, module functions, user inputs, then executable statements. Each
layer is keyed by DefId, so later layers can shadow where the language allows
without corrupting earlier registrations.

The top-level scope is also a lifetime boundary. During execution, statements
may allocate temporary arrays, circular buffers, closures, or function-local
values. Only observable bindings are copied into the output map. When the scope
ends, unobserved temporaries can be released by Python's normal reference
management. This does not make \einlang{} a native runtime, but it gives the
NumPy backend a disciplined way to avoid retaining large intermediate tensors
after a program has finished.

\textit{Concrete example.} A script-style matrix multiplication binding
executes at the top level and returns \texttt{C} as an observable output. A
program with an explicit \texttt{main} function instead resolves and calls
\texttt{main}, returning the entry-point result rather than every top-level
helper binding. Both paths register stdlib functions and builtins before
statement execution, so generated derivative code can still call helpers such
as the basis-tensor builtin.

\section{Expression Execution}

Expression execution handles literals, identifiers, arithmetic, calls, blocks,
conditionals, arrays, tuples, member access, casts, matches, where expressions,
pipelines, and builtins. Function calls push a new scope, bind parameters by
DefId, evaluate the function body, and pop the scope. Builtins include
\texttt{assert}, \texttt{print}, \texttt{len}, \texttt{typeof},
\texttt{array\_append}, \texttt{shape}, \texttt{sum}, \texttt{max},
\texttt{min}, and \texttt{\_\_basis\_tensor}.

The environment also tracks the most recent value for a DefId so escaped local
values can still be recovered in controlled cases. This is a runtime convenience
around block and function scopes; it does not change the compiler's rule that
semantic identity is DefId-based.

\section{Tensor Execution}

Lowered tensor execution chooses among vectorized, scalar, hybrid, and
specialized reduction paths using compiler-provided facts. The implementation
contains call-index analysis, recurrence analysis, vectorization helpers, and
setup utilities for temporary tensors and circular buffers. This is the largest
backend surface because it must preserve \einlang{}'s source-level tensor
semantics while exploiting NumPy when it is safe to do so.

\subsection{Lowered Tensor Execution Algorithm}

The backend executes compiler-planned lowered clauses. It may choose vectorized
or scalar code paths, but the semantic loop plan is always the lowered IR.

\begin{lstlisting}[style=pythoncode,caption={Executing a lowered Einstein value.}]
execute_lowered_einstein(lowered, binding):
    shape = evaluate compiler shape or infer from clause ranges
    output = existing output array for binding or allocate array(shape)
    env.set_value(binding.defid, output)

    recurrence_items = []
    non_recurrence_items = []
    for item in lowered.items:
        rec_dims = item.recurrence_dims_override
        if rec_dims and item.body reads binding.defid:
            recurrence_items.append(item)
            remember recurrence loops
        else:
            non_recurrence_items.append(item)

    for item in non_recurrence_items:
        execute_clause(item, output)

    if recurrence_items:
        for rec_context in execute_lowered_loops(recurrence_loops):
            for item in recurrence_items in source order:
                ok = execute_one_recurrence_step(item, rec_context, output)
                if not ok:
                    execute full clause fallback against current output
        return output

    for item in lowered.items:
        execute_clause(item, output)
    return output

execute_clause(item, output):
    strategy = item.vectorization_strategy
    if strategy is vectorized:
        evaluate body over broadcast loop axes and write slice
    else if strategy is recurrence-hybrid:
        scalarize recurrence axes and broadcast remaining axes
    else if strategy is call-scalar or scalar:
        for context in execute_lowered_loops(item.loops):
            context = execute_lowered_bindings(item.bindings, context)
            if check_lowered_guards(item.guards, context):
                output[index_tuple(context)] = evaluate(item.body, context)
\end{lstlisting}

The same helper functions execute reductions: build loop contexts, evaluate
local bindings, skip contexts whose guards fail, and combine body values using
the reduction identity and operator.

The lowered Einstein executor is organized around semantic partitions rather
than around one fixed loop implementation. Non-recurrent clauses can be
executed first because they do not depend on future writes of the same binding.
Recurrent clauses are separated so their loop order can honor recurrence
metadata. Within each clause, the vectorization strategy controls performance,
but the lowered loop plan controls meaning. This separation lets the backend
take fast NumPy paths without changing the source dependency relation.

The fallback cases are part of correctness. A recurrence step may discover that
the optimized path cannot handle a particular body, guard, or call dependency.
In that case, it executes a full clause fallback against the current output
rather than producing a partial result. Similarly, scalar and call-scalar paths
are slower but semantically direct: they enumerate loop contexts, evaluate
bindings, test guards, and write exactly one output point. The smart path is
therefore optional; the lowered IR remains the source of truth.

The recurrence partition also prevents base clauses from being repeatedly
executed inside recurrent loops. Non-recurrent clauses initialize or fill points
that do not depend on future values. Recurrent clauses then run in an order
chosen from compiler metadata. This organization is what makes boundary clauses
and interior recurrence clauses share one lowered binding without producing
incorrect overwrites. The backend does not infer this relation from names; it
receives the recurrence dimensions and source-ordered items from the compiler.

Vectorization is deliberately subordinate to guards and bindings. A clause can
be fully vectorized only when its body, local where-bindings, guards, and calls
can all be evaluated over broadcast loop axes with the same meaning as scalar
iteration. Hybrid execution scalarizes the dimensions that break that property
while preserving vector execution elsewhere. The algorithm therefore describes a
semantic ladder: vectorize when the facts prove it, hybridize when only some
axes are safe, and fall back to scalar loops when direct execution is the only
clear guarantee.

\textit{Concrete example.} A lowered matrix multiplication clause with no
guards and no loop-dependent calls can evaluate the body over broadcast arrays
and write the whole output slice through a vectorized path. A recurrence such
as \texttt{h[t] = step(h[t - 1], x[t])} is partitioned differently: the base
clause initializes \texttt{h[0]}, and the recurrent clause runs timestep by
timestep so each write sees the previous value. A guarded clause such as
\texttt{let y[i] = x[i] where x[i] > 0} may scalarize the guard check if the
facts do not prove that a vectorized mask path is safe.

\section{IREE Backend}

\pathcode{src/einlang/backends/iree.py} implements an experimental hybrid IREE
backend. It subclasses the NumPy backend and lazily compiles structurally
eligible pure functions through IREE's CPU path. The supported subset includes
scalars and tensors with simple arithmetic, casts, array literals, identifiers,
and simple block-local bindings. Unsupported functions fall back to the NumPy
execution path, so language coverage remains intact while the compiled subset
grows.

The IREE backend infers value specs from runtime arguments, emits MLIR for the
supported function body, caches compiled modules by function DefId and argument
signature, invokes compiled callables, and reports compile or runtime failures
with the original function name.

The IREE path is deliberately hybrid. A supported function can compile to IREE
while the rest of the program continues through NumPy. The cache key includes
the function DefId and argument shape/dtype signature, so the same source
function can have multiple compiled variants. Unsupported operators,
side-effectful constructs, or complex tensor control flow use the fallback path
instead of blocking execution of the whole program.

\subsection{Hybrid IREE Algorithm}

\begin{lstlisting}[style=pythoncode,caption={Lazy IREE compilation with fallback.}]
call_function_iree(func, args):
    if func is not a named Einlang function:
        return numpy_call_function(func, args)
    if body contains unsupported node or side effect:
        return numpy_call_function(func, args)

    arg_specs = infer shape and dtype from runtime args
    key = (func.defid, arg_specs.signature)
    if module_cache[key] is fallback_sentinel:
        return numpy_call_function(func, args)

    if key not in module_cache:
        preview = numpy_call_function(func, args)
        try:
            mlir = emit_mlir(func.body, arg_specs, result_spec(preview))
            module_cache[key] = compile_with_iree(mlir)
        except unsupported:
            module_cache[key] = fallback_sentinel
            return preview

    return invoke_compiled(module_cache[key], prepared_args)
\end{lstlisting}

The preview run is intentionally pragmatic: it gives the backend an actual
result shape and dtype for the supported subset while preserving correct
behavior if compilation later proves impossible.

The hybrid IREE path is designed to be opportunistic. It first rejects
functions that are not named \einlang{} functions or that contain unsupported
side effects, because compiling them would change the execution model. For
eligible functions, the backend builds a signature from runtime argument
shapes and dtypes and caches compiled modules by DefId plus signature. This
allows the same source function to compile separately for different tensor
shapes while still falling back safely when the subset is exceeded.

The preview execution is a deliberate trade-off. It may do work twice on a
successful first compile, but it gives the compiler-backed emitter a concrete
result specification and gives the runtime a correct value if MLIR emission or
IREE compilation fails. The fallback sentinel then prevents repeated failed
compilation attempts for the same signature. The result is a backend that can
grow incrementally without making unsupported programs unexecutable.

The preview run also guards result typing. MLIR emission needs to know not only
argument shapes and dtypes but also the result shape expected from the supported
function body. Rather than duplicating the full evaluator in the emitter, the
backend asks the existing NumPy path for one correct value and uses that value
as a specification. This is not ideal for peak performance, but it is a sound
engineering compromise for an experimental backend whose supported subset is
still growing.

The fallback sentinel is a negative cache. Without it, an unsupported function
called in a loop would repeatedly attempt compilation and repeatedly fail. By
storing failure for the DefId/signature pair, the backend pays the exploratory
cost once and then executes through NumPy directly. This keeps the hybrid design
predictable: adding IREE cannot make unsupported programs slower without bound,
and successful compiled paths remain isolated from functions that stay in the
interpreter.

\textit{Concrete example.} A pure helper \texttt{fn add2(x) \{ x + 2.0 \}}
called with a float matrix can be previewed once through NumPy, emitted as a
small MLIR function, and cached under the helper's DefId plus the matrix
shape/dtype signature. A helper that calls \texttt{print}, performs unsupported
control flow, or touches a complex lowered tensor node receives a fallback
sentinel for that signature and continues through NumPy. Repeated calls to the
unsupported helper do not repeatedly attempt compilation.

\section{Memory Model}

\einlang{} does not implement a custom garbage collector. Ordinary values are
Python objects and NumPy arrays. The runtime reduces memory pressure by clearing
scope dictionaries on exit, avoiding state reuse between executions, using
circular buffers for bounded recurrences, and materializing lazy values only
when observation requires it. These choices are runtime engineering decisions,
not a replacement for a native heap or precise GC.

\subsection{Scope Lifetime Algorithm}

\begin{lstlisting}[style=pythoncode,caption={DefId-keyed scope lifetime.}]
with env.scope():
    push empty dict onto scope stack
    execute block or function body
    on exit:
        copy escaped or output values already recorded by DefId
        clear popped scope dictionary
        pop scope

get_value(defid):
    for scope in reversed(scope_stack):
        if defid in scope:
            return scope[defid]
    return last_values.get(defid) for controlled escaped-local cases
\end{lstlisting}

This is not memory management in the native-runtime sense, but it matters for
large NumPy arrays. A completed block or function should not keep temporary
tensors alive through a stale scope dictionary.

The scope algorithm is small because Python and NumPy own actual allocation,
but it still enforces an important lifetime discipline. Every block or function
pushes a dictionary keyed by DefId; leaving that block clears the dictionary so
temporary arrays are not retained by the runtime environment. Values that must
escape are copied to output storage or remembered in the controlled
\texttt{last\_values} path before the scope is popped. This gives the runtime a
predictable ownership boundary without implementing a custom heap.

Lookup walks from inner to outer scopes, matching lexical shadowing after name
resolution. Because every binding has a DefId, two variables with the same
printed name never conflict in the environment. This is the runtime mirror of
the resolver invariant: user-facing names are useful for diagnostics and API
inputs, but executable storage is keyed by semantic identity.

The \texttt{last\_values} path is intentionally narrow. It supports controlled
cases where a value must be recoverable after its lexical scope is gone, but it
does not turn the environment into a global variable table. Ordinary lookup
still respects the active scope stack, and source programs cannot address local
values by string after resolution. This keeps the runtime convenient for
execution bookkeeping without weakening lexical scoping.

Clearing dictionaries on scope exit is especially relevant for tensor programs.
A temporary in a scalar example is cheap; a temporary in a convolution,
recurrence, or generated Jacobian can be large enough to dominate memory use.
Because Python frees arrays only when references disappear, stale scope entries
would keep tensors alive even after they are semantically dead. The scope
algorithm is therefore a small piece of runtime hygiene that supports the
larger compiler goal of making storage decisions explicit.

\textit{Concrete example.} In a norm helper, a temporary tensor \texttt{sq}
may store \texttt{x[i] * x[i]} while the function body computes
\texttt{sum[i](sq[i])}. The temporary is stored in the function scope while the
body runs. After the return value is computed, the scope dictionary holding
\texttt{sq} is cleared, so the NumPy array can be released when no other
references remain. If the returned value is bound at the top level, that value
is copied to output storage by DefId before the function scope disappears.

\section{Debugging and Profiling}

The CLI and backend expose practical debugging tools: IR dumps, statement
profiling, line-bucket profiling, function profiling, block profiling,
reduction-path profiling, vectorization summaries, recurrence-block debug
output, and structured autodiff debug logs. These hooks are implementation
features, not just developer conveniences: they make it possible to validate
whether a large example is exercising the intended compiler and backend paths.
